\section{Theoretische Grundlagen}
\label{sec:grundlagen}

Dieses Kapitel legt die mathematischen, physikalischen und informationstechnischen Grundlagen dar, auf denen die Simulationspipeline aufbaut.
Das Verständnis dieser Konzepte ist Voraussetzung für die Bewertung der in späteren Kapiteln beschriebenen Implementierungsentscheidungen.

\subsection{Akustische Grundlagen}
\label{subsec:akustik}

\paragraph{Das Inverse Quadratgesetz}

Im freien, reflexionsfreien Schallfeld breitet sich Schall kugelförmig vom Emitter aus.
Da die abgestrahlte Schallleistung $P$ über eine mit $r^2$ wachsende Kugeloberfläche
$A = 4\pi r^2$ verteilt wird, nimmt die Schallintensität $I$ mit dem Quadrat der Entfernung ab:
\[I(r) = \frac{P}{4\pi r^2}\]
Da die Schallintensität proportional zum Quadrat des Schalldrucks ist ($I \propto p^2$), folgt für den Schalldruck:
\[p(r) \propto \frac{1}{r}\]
Eine Verdopplung des Abstands zur Quelle halbiert den Schalldruck, was einer Pegelminderung von $6\,\text{dB}$ entspricht.
In der Simulation wird dieser geometrische Energieverlust implizit berücksichtigt:
Strahlen, die einen längeren Weg zur Mikrofonkapsel zurücklegen, treffen aufgrund der divergenten Geometrie mit einem entsprechend geringeren Druckbeitrag ein.

\paragraph{Materialabsorption}

Trifft ein Schallstrahl auf eine Raumgrenze, wird ein Teil der Schallenergie absorbiert.
Dieser Vorgang wird durch den dimensionslosen Absorptionskoeffizienten $\alpha \in [0{,}1]$ beschrieben. Der Restdruck nach einer einzelnen Reflexion berechnet sich zu:
\[p_{\text{ref}} = p_{\text{ein}} \cdot \sqrt{1 - \alpha}\]
Bei einer Sequenz von $n$ Reflexionen mit den jeweiligen Koeffizienten $\alpha_1, \ldots, \alpha_n$ gilt:
\[p_n = p_0 \cdot \prod_{k=1}^{n}\sqrt{1 - \alpha_k}\]
Ein vollständig absorbierendes Material ($\alpha = 1{,}0$) vernichtet den Strahl bei der ersten Reflexion, während ein perfekt reflektierendes Material ($\alpha = 0{,}0$) keine Energie entzieht. Die Simulation terminiert Strahlen, sobald ihr verbleibender Druckwert den konfigurierbaren Grenzwert \texttt{min\_pressure} unterschreitet.

\paragraph{Frühe Reflexionen und späte Nachhallzeit}

Die akustische Impulsantwort eines Raumes lässt sich zeitlich in zwei charakteristische Phasen unterteilen.
Innerhalb der ersten ca. $80\,\text{ms}$ nach dem Direktschall treffen die \emph{frühen Reflexionen} (Early Reflections) ein: diskrete, noch einzeln wahrnehmbare Echos über klar definierte, kurze Ausbreitungspfade, die das Raumgefühl und die Tiefenwahrnehmung maßgeblich prägen.
Mit zunehmender Zeit überlagern sich Reflexionen höherer Ordnung zu einem dichten, statistisch gleichverteilten Nachhallfeld. Diese \emph{späte Nachhallzeit} (Late Reverberation) klingt exponentiell ab und wird durch die Nachhallzeit $RT_{60}$ charakterisiert -- der Zeitraum, in dem der Schalldruckpegel um $60\,\text{dB}$ abfällt.
Das stochastische Raytracing erfasst beide Phasen natürlich, da sowohl kurze Direktpfade als auch komplexe Mehrfachreflexionen in der Impulsantwort akkumuliert werden.


\subsection{Monte-Carlo-Integration und Raytracing}
\label{subsec:montecarlo}

\paragraph{Das akustische Integrationsproblem}

Der Schalldruck am Empfänger ist das Integral über alle möglichen Schallpfade vom Emitter, gewichtet nach Weglänge und materialabhängigen Verlusten. Da die Anzahl dieser Pfade durch Mehrfachreflexionen kombinatorisch explodiert, ist eine analytische Auswertung für beliebige Raumgeometrien nicht praktikabel. Deterministisches Raytracing kann zwar exakte Pfade verfolgen, skaliert jedoch exponentiell mit der Anzahl der Reflexionen.

\paragraph{Der Monte-Carlo-Ansatz}

Die Monte-Carlo-Integration umgeht dieses Problem durch stochastische Stichprobennahme.
Statt alle Pfade zu berechnen, schätzt man den Erwartungswert eines Integrals durch den arithmetischen Mittelwert von $N$ zufällig gezogenen Stichproben:
\[\mathbb{E}[f(X)] \approx \frac{1}{N} \sum_{i=1}^{N} f(x_i)\]
Nach dem Gesetz der großen Zahlen konvergiert dieser Schätzer für $N \to \infty$ gegen den wahren Integralwert. Der statistische Fehler nimmt dabei mit $\frac{1}{\sqrt{N}}$ ab: ein viermal so großes Stichprobenvolumen halbiert den Fehler.

\paragraph{Anwendung: Stochastisches Raytracing}

In der Simulation entspricht eine Stichprobe einem Schallstrahl, der in einer zufälligen Richtung $\theta \in [0, 2\pi)$ vom Emitter ausgesendet wird:
\[\vec{D} = (\cos\theta,\; \sin\theta), \quad \theta \sim \mathcal{U}(0,\; 2\pi)\]
Jeder Strahl verfolgt seinen Ausbreitungspfad, reflektiert an Wänden und registriert einen Treffer, sobald er die Mikrofonkapsel schneidet. Der akkumulierte Beitrag dieses Treffers -- Laufzeit $\tau$ und Restdruck $p$ -- repräsentiert eine Stichprobe des akustischen Integrals. Über $N = 10^6$ bis $10^8$ Strahlen konvergiert die Gesamtheit der Treffer zur vollständigen Impulsantwort des Raumes.
Der wesentliche Vorteil gegenüber deterministischen Verfahren liegt in der trivialen Parallelisierbarkeit: Da jeder Strahl unabhängig von allen anderen ist, lassen sich die $N$ Stichproben direkt auf beliebig viele CPU-Kerne aufteilen (siehe Abschnitt~\ref{subsec:amdahl}).


\subsection{Digitale Signalverarbeitung: Impulsantwort und schnelle Faltung}
\label{subsec:dsp}

\paragraph{Die Impulsantwort (IR)}

Ein lineares, zeitinvariantes System -- wie ein Raum aus akustischer Perspektive -- ist vollständig durch seine \emph{Impulsantwort} $h(t)$ charakterisiert: die Ausgabe des Systems bei einem idealen Dirac-Impuls $\delta(t)$ am Eingang.
Die von der Rust-Engine berechnete Menge an (Laufzeit, Druck)-Paaren $\{(\tau_i, p_i)\}_{i=1}^{M}$ stellt eine diskrete Approximation dieser Impulsantwort dar:
\[h(t) \approx \sum_{i=1}^{M} p_i \cdot \delta(t - \tau_i)\]
Das Python-Modul diskretisiert diese Darstellung in ein sample-genaues Array, indem jede Verzögerung $\tau_i$ in den Sample-Index $n_i = \lfloor \tau_i \cdot f_s \rceil$ umgerechnet und Druckwerte akkumuliert werden (bei $f_s = 44100\,\text{Hz}$).

\paragraph{Das Faltungsintegral}

Das Ausgangssignal $y(t)$ -- das hallbehaftete Audio -- ergibt sich aus der linearen Faltung des Eingangssignals $x(t)$ mit der Impulsantwort $h(t)$:
\[y(t) = (x * h)(t) = \int_{-\infty}^{\infty} x(\tau)\, h(t - \tau)\, d\tau\]
In der diskreten Signalverarbeitung wird dieses Integral zur Summe:
\[y[n] = \sum_{k=0}^{M-1} h[k] \cdot x[n - k]\]
Die direkte Auswertung hat eine Rechenkomplexität von $\mathcal{O}(N \cdot M)$, wobei $N$ die Länge des Audiosignals und $M$ die Länge der Impulsantwort bezeichnet. Bei realen Impulsantworten mit $M$ von mehreren hunderttausend Samples ist eine direkte Berechnung unpraktikabel.

\paragraph{Schnelle Faltung via FFT (Overlap-Add)}

Das \emph{Faltungstheorem} der Fouriertransformation besagt, dass eine Faltung im Zeitbereich einem punktweisen Produkt im Frequenzbereich entspricht:
\[Y(f) = \mathcal{F}\{x\}(f) \cdot \mathcal{F}\{h\}(f) \implies y[n] = \mathcal{F}^{-1}\{Y(f)\}\]
Durch den Einsatz der Fast Fourier Transform (FFT) sinkt die Komplexität auf $\mathcal{O}((N+M)\log(N+M))$, was bei großen $M$ einen massiven Speedup bedeutet.

Da das Audiosignal wesentlich länger sein kann als eine einzelne FFT erlaubt, wird das \emph{Overlap-Add}-Verfahren eingesetzt: Das Signal $x[n]$ wird in überlappungsfreie Blöcke der Länge $L$ aufgeteilt, jeder Block wird separat mit $h[n]$ gefaltet, und die resultierenden Ausgabeblöcke -- je $L + M - 1$ Samples lang -- werden additiv überlagert. Die Implementierung nutzt hierfür \texttt{scipy.signal.oaconvolve}, welche diesen Algorithmus effizient realisiert.


\subsection{Theoretische Grundlagen paralleler Systeme}
\label{subsec:amdahl}

Um die in späteren Kapiteln (siehe Kapitel~\ref{sec:evaluation}) gemessenen Laufzeiten einordnen zu können, wird der theoretische Rahmen zur Bewertung paralleler Systeme herangezogen.

\paragraph{Speedup}
Der Speedup $S(n)$ bei Verwendung von $n$ parallelen Verarbeitungseinheiten (hier: Nodes oder CPU-Threads) ist definiert als das Verhältnis der seriellen Laufzeit $T(1)$ zur parallelen Laufzeit $T(n)$:
\begin{equation}
    S(n) = \frac{T(1)}{T(n)}
    \label{eq:speedup}
\end{equation}
Ein Speedup von $S(n) = n$ wird als \emph{linear} bezeichnet und stellt den theoretischen Idealfall dar, bei dem eine Verdopplung der Ressourcen die Laufzeit exakt halbiert.

\paragraph{Amdahls Gesetz}
In der Praxis lässt sich jedoch nicht jeder Programmteil parallelisieren. Amdahls Gesetz beschreibt den maximal erreichbaren Speedup in Abhängigkeit vom parallelisierbaren Anteil $p$ eines Programms (bei einem seriellen, nicht parallelisierbaren Restanteil von $1-p$):
\begin{equation}
    S(n) = \frac{1}{(1-p) + \dfrac{p}{n}}
    \label{eq:amdahl}
\end{equation}
Für $n \to \infty$ konvergiert der Speedup gegen den Grenzwert $\frac{1}{1-p}$. Das bedeutet: Selbst bei einer beliebig großen Anzahl an parallelisierten Threads oder Nodes limitiert bereits ein kleiner serieller Anteil den maximal erreichbaren Gesamtspeedup. Ein Programm mit beispielsweise $10\,\%$ seriellem Anteil ($p = 0{,}9$) kann demnach nie schneller als das 10-fache der Ausgangslaufzeit werden, unabhängig von der Anzahl eingesetzter Recheneinheiten.

\paragraph{Anwendung auf das Ray-Tracing-Problem}
Das hier vorliegende stochastische Simulationsproblem gehört zur Klasse der \emph{embarrassingly parallel} Probleme: Jeder der $N_{\text{total}}$ Schallstrahlen kann vollständig unabhängig von allen anderen berechnet werden, da während der Ray-Generierung kein Datenaustausch zwischen den Berechnungsinstanzen erforderlich ist. Der serielle Anteil $1-p$ beschränkt sich im Systemaufbau auf die Orchestrierungsphasen (Verbindungsaufbau, Dateitransfer via \texttt{scp} bei Clustern, sowie die Zusammenführung der Ergebnisse). Dieser Anteil ist im Vergleich zur Berechnungszeit minimal, weshalb ein nahezu linearer Speedup und eine lineare Durchsatzsteigerung erwartet und empirisch belegt werden kann.